% BHU PYQ -- Manually transcribed & error-corrected
% Paper: GLB-604: Economic Geology
% Exam: B.Sc. (Hons.) Semester VI Examination, 2013-14

\documentclass[11pt,a4paper]{article}
\usepackage[a4paper,top=2.5cm,bottom=2.5cm,left=2.8cm,right=2.8cm,headheight=14pt]{geometry}
\usepackage{amsmath,amssymb}
\usepackage[T1]{fontenc}
\usepackage[utf8]{inputenc}
\usepackage{lmodern,microtype,fancyhdr,enumitem,hyperref,lastpage,parskip}

\hypersetup{colorlinks=true,urlcolor=black,linkcolor=black,
  pdftitle={GLB-604 Economic Geology -- BHU B.Sc. Sem VI 2013-14}}
\pagestyle{fancy}\fancyhf{}
\renewcommand{\headrulewidth}{0.4pt}\renewcommand{\footrulewidth}{0.4pt}
\fancyhead[L]{\small\textit{Banaras Hindu University}}
\fancyhead[R]{\small\textit{GLB-604: Economic Geology}}
\fancyfoot[C]{\small Page \thepage\ of \pageref{LastPage}}

\newlist{parts}{enumerate}{2}
\setlist[parts,1]{label=(\alph*),leftmargin=*,itemsep=5pt,topsep=3pt}
\newcommand{\pts}[1]{\hfill\textit{[#1]}}

\begin{document}
\begin{center}
  {\large\bfseries BANARAS HINDU UNIVERSITY}\\[2pt]
  {\normalsize B.Sc. (Hons.) --- Semester VI Examination, 2013-14}\\[6pt]
  \rule{0.8\linewidth}{0.5pt}\\[6pt]
  {\large\bfseries Paper: GLB-604 --- Economic Geology}\\[4pt]
  {\normalsize Time: 3 Hours \hfill Full Marks: 70}\\[2pt]
  {\small Ref. No.: BS/Sem VI/228}
\end{center}
\rule{\linewidth}{0.4pt}
\smallskip
\noindent\textit{Note: The figures in the right-hand margin indicate marks. Answer \textbf{five} questions in all, selecting \textbf{two} each from Sections A and B. Section-C is compulsory.}
\bigskip

\bigskip\noindent{\large\bfseries SECTION --- A}
\rule{\linewidth}{0.4pt}\smallskip

\noindent\textbf{Question 1.} Discuss the origin of petroleum and also write its important distribution in India. \pts{14}

\medskip\noindent\textbf{Question 2.} Write an essay on hydrothermal process of formation of ore deposit with Indian example. \pts{14}

\medskip\noindent\textbf{Question 3.} What is placer deposit? Discuss the mechanical concentration process of formation of ore deposit with Indian example. \pts{14}

\medskip\noindent\textbf{Question 4.} Enlist the important minerals/rocks raw material used in refractory industry. Write their chemical composition, specification (if any) and important distribution in India. \pts{14}

\bigskip\noindent{\large\bfseries SECTION --- B}
\rule{\linewidth}{0.4pt}\smallskip

\noindent\textbf{Question 5.} Give the chemical compositions, physical properties, specification for different industrial uses and important distribution in India of \textbf{any two} of the following: \pts{$2 \times 7 = 14$}
\begin{parts}
  \item Halite
  \item Phosphorite
  \item Pyrite
\end{parts}

\medskip\noindent\textbf{Question 6.} Write note on \textbf{any two} of the following: \pts{$2 \times 7 = 14$}
\begin{parts}
  \item Very common forms and structures of ore deposits
  \item Sedimentation process
  \item Contact metasomatic process
\end{parts}

\medskip\noindent\textbf{Question 7.} Answer the following: \pts{$2 \times 7 = 14$}
\begin{parts}
  \item Discuss the important distribution of coal (bituminous and lignite) in India.
  \item Write a note on building stone.
\end{parts}

\bigskip\noindent{\large\bfseries SECTION --- C}
\rule{\linewidth}{0.4pt}\smallskip

\noindent\textbf{Question 8.} Answer the following: \pts{14}
\begin{parts}
  \item Give neat sketch of supergene sulphide enrichment process. \pts{4}
  \item Differentiate between contact metamorphism and contact metasomatism. \pts{2}
  \item Name the host rock associated with---
  \begin{enumerate}[label=(\roman*),leftmargin=*,itemsep=2pt]
    \item Amorphous variety of magnesite
    \item Chromite deposit
    \item Diamond
    \item Mica deposit
  \end{enumerate} \pts{$4 \times \frac{1}{2} = 2$}
  \item Write in brief the different stages of coal formation. \pts{3}
  \item Write in brief the basic relation between paragenesis and zoning. \pts{3}
\end{parts}

\vfill\rule{\linewidth}{0.4pt}
\begin{center}\small
\href{https://bhu-exam.vercel.app/}{Click here to visit our website}\\[4pt]
\href{https://me-aryan.vercel.app/}{Click here to visit our contributor}\\[4pt]
\href{https://quashan-me.vercel.app/}{Click here to visit our contributor}
\end{center}
\end{document}
